% Drawing 138/07-9 -- folder 138, batch 07, archive page 126 (PDF page 127).
% Transcription: un segment $[0,1]$ avec son milieu marqué ; un pentagone aux sommets noircis et aux milieux d'arêtes marqués d'un trait.
% Kind: segment with midpoint; pentagon. Method 2 (schematic redraw from the facsimile). Status: ok.
% Notes: The segment [0, 1] with its midpoint marked; below, a second short segment with a point in the middle (by « f(z) - 1/2 »); a pentagon with blackened vertices and edge midpoints marked by a stroke (one side drawn with extra dots). Formulas around are text.
% Coordinates: 1 unit = 1 mm = 1/100 of the crop width (bbox in page fractions: [0.315, 0.615, 0.584, 0.734]).
% Redrawn by Claude (claude-opus-5-5), issue #73, 2026-11-07; not checked by a human; not compiled (no TeX in the session).
% Rights: reproduces part of the work -- behind RENDER_TRANSCRIPTS, never published by itself.
\documentclass[tikz,border=3pt]{standalone}
\usepackage{amsmath,amssymb}
\usetikzlibrary{patterns,arrows.meta}
\tikzset{pencil/.style={black!45}, struck/.style={black!35}, mark/.style={circle,draw,inner sep=1.3pt}, pt/.style={circle,fill,inner sep=0.9pt}, opt/.style={circle,draw,fill=white,inner sep=1.1pt}}
\begin{document}
\begin{tikzpicture}[x=1mm,y=1mm,line cap=round,line join=round,font=\scriptsize]
\draw[very thick] (53.1,-18.8) -- (81.2,-17.2);
\draw (67.2,-15) -- (67.2,-20.6);
\node[inner sep=1pt] at (53.1,-10.9) {\tiny $0$};
\node[inner sep=1pt] at (81.2,-9.4) {\tiny $1$};
\draw (6.2,-43.8) -- (34.4,-44.4);
\node[pt] at (18.8,-44.1) {};
\draw (53.1,-34.4) -- (75,-28.1) -- (81.2,-46.9) -- (62.5,-56.2) -- (51.6,-46.9) -- cycle;
\node[pt] at (53.1,-34.4) {};
\node[pt] at (75,-28.1) {};
\node[pt] at (81.2,-46.9) {};
\node[pt] at (62.5,-56.2) {};
\node[pt] at (51.6,-46.9) {};
\node[pt] at (59.4,-32.8) {};
\node[pt] at (67.2,-31.2) {};
\draw (75.6,-35) -- (80.6,-40);
\draw (69.4,-49.1) -- (74.4,-54.1);
\draw (53.8,-49.1) -- (58.8,-54.1);
\draw (49.7,-38.1) -- (54.7,-43.1);
\end{tikzpicture}
\end{document}
